\chapter*{Аннотация}       % или \chapter*{Аннотация}, если хотите без номера
Работа посвящена исследованию спектральных характеристик внутренних представлений моделей компьютерного зрения: свёрточных сетей ResNet, визуальных трансформеров ViT и иерархических трансформеров Swin. Цель исследования — количественно описать, как по мере углубления модели меняется распределение энергии представлений по пространственным частотам, и сравнить динамику между CNN и Transformer архитектурами.

В работе реализован единый протокол спектрального анализа. Для ResNet анализировались выходы основных блоков, для ViT и Swin последовательности токенов преобразовывались в двумерную решётку патчей. Для выбранных уровней выполнялись вычитание пространственного среднего, ДПФ, спектр мощности, усреднение по батчу и каналам и радиальное усреднение. Для сравнения использовались спектральный центроид, доля низких частот и скорость сжатия.

Эксперименты показали, что ResNet смещает представления к низким частотам: центроид уменьшается, а доля низкочастотной энергии возрастает. ViT имеет немонотонную динамику: в ряде моделей центроид на средних и поздних блоках не снижается, а возрастает. Swin занимает промежуточное положение: его иерархия приводит к низкочастотному финальному представлению, но сжатие спектра менее резкое, чем в CNN.

Результаты показывают, что тип архитектуры сильнее влияет на спектральную динамику, чем глубина внутри одного семейства. Подход может использоваться для диагностики моделей, особенно в задачах, чувствительных к деталям и локальным структурам.